%\section{Senderausgangsleistung}
%\label{section:senderausgangsleistung}

\begin{frame}
    \frametitle{Messung der Senderausgangsleistung}
    \begin{columns}
        \begin{column}{0.48\linewidth}
            \begin{figure}
                \centering
                \DARCimage{1.0\linewidth}{916include}
                \caption{Messung der Senderausgangsleistung}
            \end{figure}
        \end{column}
        \begin{column}{0.48\linewidth}
            \begin{itemize}
                \item Direkt am Senderausgang
                \item Ohne Zusatzgeräte, Filter oder Kabel
                \item Bei SSB $\rightarrow$ mit Modulation
                \item Ein- oder Zweitonaussteuerung, aber keine Sprache
                \item Messung der maximalen \emph{Hüllkurvenleistung} (PEP)
                \item Spitzenleistung des Senders bei maximaler Aussteuerung
            \end{itemize}
        \end{column}
    \end{columns}
\end{frame}

\begin{frame}
    \only<1>{
        \begin{QQuestion}{EF401}{Die Ausgangsleistung eines Senders ist die unmittelbar nach~...}
            {dem Senderausgang gemessene Summe aus vorlaufender und rücklaufender Leistung.}
            {dem Senderausgang gemessene Differenz aus vorlaufender und rücklaufender Leistung.}
            {der Antenne messbaren Leistung, die durch ein Feldstärkenmessgerät im Nahfeld ermittelt werden kann.}
            {dem Senderausgang messbare Leistung, bevor sie Zusatzgeräte durchläuft.}
        \end{QQuestion}

    }
    \only<2>{
        \begin{QQuestion}{EF401}{Die Ausgangsleistung eines Senders ist die unmittelbar nach~...}
            {dem Senderausgang gemessene Summe aus vorlaufender und rücklaufender Leistung.}
            {dem Senderausgang gemessene Differenz aus vorlaufender und rücklaufender Leistung.}
            {der Antenne messbaren Leistung, die durch ein Feldstärkenmessgerät im Nahfeld ermittelt werden kann.}
            {\textbf{\textcolor{DARCgreen}{dem Senderausgang messbare Leistung, bevor sie Zusatzgeräte durchläuft.}}}
        \end{QQuestion}

    }
\end{frame}

\begin{frame}
    \only<1>{
        \begin{QQuestion}{EF402}{Wie und wo wird die Ausgangsleistung eines SSB-Senders gemessen? Die maximale Hüllkurvenleistung (PEP) wird gemessen...}
            {zwischen Antennentuner und Speisepunkt bei Sprachmodulation.}
            {zwischen Antennentuner und Speisepunkt der Antenne mit unmoduliertem Träger.}
            {direkt am Senderausgang bei Ein- oder Zweitonaussteuerung.}
            {direkt am Senderausgang mit unmoduliertem Träger.}
        \end{QQuestion}

    }
    \only<2>{
        \begin{QQuestion}{EF402}{Wie und wo wird die Ausgangsleistung eines SSB-Senders gemessen? Die maximale Hüllkurvenleistung (PEP) wird gemessen...}
            {zwischen Antennentuner und Speisepunkt bei Sprachmodulation.}
            {zwischen Antennentuner und Speisepunkt der Antenne mit unmoduliertem Träger.}
            {\textbf{\textcolor{DARCgreen}{direkt am Senderausgang bei Ein- oder Zweitonaussteuerung.}}}
            {direkt am Senderausgang mit unmoduliertem Träger.}
        \end{QQuestion}

    }
\end{frame}%ENDCONTENT
